\documentclass[conference]{IEEEtran}

\usepackage{fontspec}
\usepackage{xeCJK}
\usepackage{hyperref}
\usepackage{booktabs}
\usepackage{url}
\usepackage{listings}
\usepackage{graphicx}
\usepackage{xcolor}

\setCJKmainfont{Hiragino Mincho ProN}
\setCJKsansfont{Hiragino Sans}
\setCJKmonofont{Hiragino Sans}
\XeTeXlinebreaklocale "ja"
\XeTeXlinebreakskip = 0pt plus 1pt

\hypersetup{
  colorlinks=true,
  linkcolor=black,
  urlcolor=blue,
  citecolor=black
}

\lstset{
  basicstyle=\small\ttfamily,
  breaklines=true,
  frame=single,
  columns=fullflexible
}

\begin{document}

\title{Associative Context Memory: コンテキストウィンドウ依存を超えるLLMコーディングエージェントのための永続的体験ベースメモリ}

\author{
  \IEEEauthorblockN{Hiroshi Yamato}
  \IEEEauthorblockA{
    Keio University / Signal compose Inc.\\
    yamato@signalcompose.com
  }
}

\maketitle

\begin{abstract}
大規模言語モデル（LLM）ベースのコーディングエージェントは、根本的な制約を抱えている。その記憶はコンテキストウィンドウに完全に束縛されており、トークン数の増加に伴い品質が劣化する——Context Rotとして知られる現象である。コンテキスト圧縮、セッション分割、検索拡張生成（RAG）といった既存の緩和策は症状に対処するものの、コンテキストウィンドウサイズへの根本的な依存を解消するものではない。本研究では、\textbf{Associative Context Memory (ACM)}を提案する。ACMはLLMコーディングエージェントのための外部永続メモリアーキテクチャであり、暗黙的なユーザーフィードバックシグナルから成功・失敗体験を捕捉し、コンテキストウィンドウ容量に依存することなくセッションをまたいだ学習を可能にする。本研究の主要な貢献は、コーディングエージェントインタラクションにおける暗黙的フィードバックシグナルの分類体系とその強度序列である。ACMをMCP（Model Context Protocol）サーバーとして設計し、既存のLLMコーディングエージェントエコシステムとの互換性を確保する。本論文はACMをポジションペーパーおよびシステム提案として提示するものであり、Phase~0実現可能性プローブを含む。
\end{abstract}

\section{はじめに}

Claude CodeやOpenAI Codexに代表されるLLMベースのコーディングエージェントは、ソフトウェア開発タスクにおいて顕著な能力を示している。しかし、これらは根本的なアーキテクチャ上の制約を共有している。すなわち、すべてのコンテキスト——会話履歴、ツール出力、システムプロンプト、検索されたドキュメント——が有限のコンテキストウィンドウに収まる必要がある。セッション内でトークン数が増加するにつれ、モデルの性能は測定可能な形で劣化する。近年の実証研究はこの現象を\textbf{Context Rot} \cite{ChromaResearch2025}として体系的に特徴づけ、トークン全体に注意が分散し、コンテキストが蓄積されるにつれてタスク成功確率が低下することを示している。

この問題に対する現在のアプローチは2つのカテゴリに分類される。

\textbf{圧縮・分割}アプローチ（auto-compact、Plan Mode、手動セッションリセット）は、Context Rotをセッション管理の問題として扱う。コンテキストサイズの削減により劣化を緩和するが、セッションをまたいで知識を引き継ぐ機構を持たない。新しいセッションはすべてゼロから始まる——我々はこれを\textbf{「健忘エージェント（amnesiac agent）問題」}と呼ぶ。

\textbf{検索拡張}アプローチ（Serena \cite{OraiosAI2025Serena}、RAGベースのコンテキスト拡張）は、エージェントの関連情報へのアクセスを改善するが、何が成功し何が失敗したかという経験的知識の捕捉は行わない。

いずれのカテゴリも、より深い問いに答えていない。\textbf{エージェントは自身のインタラクション履歴から、コンテキストウィンドウとセッションをまたいで持続する形で学習できるか？}

人間の認知は示唆に富む対照を提供する。経験豊富な開発者は過去のデバッグセッションの逐語的な記憶を保持しない。代わりに\emph{連想パターン}を蓄積する。この知識は疎で、圧縮されており、連想検索の手がかりによってインデックスされている。

本研究では、LLMコーディングエージェントのための永続的な外部メモリ層である\textbf{Associative Context Memory (ACM)}を提案する。ACMは以下を実現する。

\begin{enumerate}
\item 暗黙的なユーザーフィードバックシグナルから成功・失敗体験を捕捉する
\item 連想キーでインデックスされた構造化・検索可能な形式で体験を保存する
\item 過度なコンテキスト消費なしに、関連する過去の体験を新しいセッションに注入する
\item 持続的な性能のためのコンテキストウィンドウサイズへの依存を低減する
\end{enumerate}

本研究の中心的な新規性は、コーディングエージェントインタラクションにおいて利用可能な\textbf{暗黙的フィードバックシグナルの分類体系}である。我々は、インタラプトイベント——Ctrl+Cまたは同等の手段によるエージェント実行の強制停止——を、特にその直後の対話と組み合わせた場合に、ユーザー不満の最も強力なシグナルとして同定する。このシグナルはPTY（擬似端末）レベルで容易に観測可能であるにもかかわらず、先行研究では見過ごされてきた。

\section{背景と関連研究}

\subsection{コンテキストウィンドウの制約とContext Rot}

Transformerアーキテクチャのself-attention機構はコンテキストウィンドウ内のすべてのトークンに対して動作し、計算量はトークン数$n$に対して$O(n^2)$でスケーリングする。Chroma Research \cite{ChromaResearch2025}はこれを「Context Rot」として体系的に特徴づけ、Claude Opus~4やSonnet~4を含む最先端モデルであってもコンテキストが充填されるにつれて測定可能な性能劣化を示すことを実証した。

実務者の間では、コンテキストウィンドウの最初の約40\%（「Smart zone」）の範囲内で性能が最も信頼できるという経験則が観察されている \cite{Huntley2025}。

\subsection{コーディングエージェントにおけるコンテキスト管理}

\textbf{Auto-compact}: コンテキストが容量に近づいた際のセッション履歴の要約。重大な弱点として、要約がコンテキスト容量の限界付近で動作する劣化したモデルによって実行される。

\textbf{Plan Mode (Claude Code)}: 計画と実装を別個のセッションに分離する。効果的だが、手動のワークフロー規律を必要とし、クロスセッションメモリを提供しない。

\textbf{CLAUDE.mdと永続的な指示} \cite{Anthropic2025}: セッション開始時に読み込まれる静的ドキュメント。安定したプロジェクトコンテキストには有用だが、インタラクション履歴に適応しない。

\subsection{エージェントメモリシステム}

LLMエージェントのための永続メモリの研究領域は急速に発展している。

\textbf{MemGPT / Letta} \cite{Packer2023MemGPT}は、コンテキストウィンドウを「RAM」、外部ストレージを「ディスク」として扱う階層型メモリアーキテクチャを導入した。最も影響力のある汎用エージェントメモリアーキテクチャであるが、結果に基づく重み付けメモリ選択やコーディング固有のフィードバックシグナルには対応していない。

\textbf{Mem0およびA-MEM} \cite{Chhikara2025Mem0}は、セマンティック検索とグラフベースの組織化を備えたメモリ層を提供する。A-MEM \cite{Xu2025AMEM}は動的なノート構築とリンク生成を通じて柔軟性を向上させている。いずれのシステムもメモリエントリを結果の品質によって区別しない。

\textbf{A-MAC} \cite{Zhang2026AMAC}は、ACMに構造的に最も類似した先行研究である。メモリ受入を5つの価値シグナルに分解するが、コーディングエージェントインタラクションからの行動シグナルを受入基準として利用しない。

\textbf{REMEMBERER} \cite{Zhang2023REMEMBERER}は、Q値を含むインタラクション記録としてエピソディックメモリを実装し、ベースLLMのファインチューニングなしに強化学習によって更新する。

\textbf{Self-Generated In-Context Examples} \cite{Sarukkai2025SelfGenICE}は、成功したタスク軌跡をin-context examplesとして保存し、ALFWorldで73\%$\rightarrow$89\%の性能向上を実証している。ACMはこの方向を拡張し、(1) 失敗体験を含め、(2) シグナル強度分類体系を導入し、(3) コーディングエージェントに特化する。

\textbf{Letta Code} \cite{LettaCode2025}は、MemGPT/Lettaアーキテクチャに基づくメモリファーストのコーディングエージェントである。Terminal-Benchで優れた結果を達成しているが、PTYレベルの行動シグナルをメモリ品質指標として活用していない。

\textbf{How Memory Management Impacts LLM Agents} \cite{Xiong2025MemoryMgmt}は、LLMエージェントの「experience-following property」——エージェントが誤ったものを含むメモリからのパターンを複製する傾向——を実証的に示している。この知見はACMにおける成功・失敗エントリの区別を直接的に動機づける。

\textbf{Reflexion} \cite{Shinn2023Reflexion}は、タスク試行の失敗後に言語的な自己フィードバックを生成し、エピソディックメモリとして保存する。重要な違いとして、Reflexionはエージェント自身の自己省察によってフィードバックを生成するのに対し、ACMはユーザーからの暗黙的行動シグナルを捕捉する。また、Reflexionはメモリを検索可能な外部ストアとしてセッション間で永続化しない。\footnote{Reflexionは単一タスク内のリトライ間でエピソディックメモリを永続化するが、異なるタスクやコードベース間で利用可能な汎用的検索可能ストアを提供しない。}

\textbf{ExpeL} \cite{Zhao2023ExpeL}は、過去のタスク軌跡から転移可能な知識を抽出する。ACMはリアルタイムの暗黙的フィードバックシグナルに焦点を当てる点で異なる。ExPeLの回顧的軌跡分析にはレイテンシが伴い、知識抽出はタスク完了後に行われるため、セッション中には利用できない。ACMのリアルタイムシグナル捕捉は、将来の拡張においてセッション内の行動調整を支援し得る。

\textbf{Voyager} \cite{Wang2023Voyager}は、Minecraft環境における成功したタスク実行から永続的なスキルライブラリを構築する。ACMはこの方向をコーディングエージェントインタラクションに拡張し、暗黙的ユーザーフィードバックを使用して失敗体験も追加する。

すべての先行メモリシステムに共通する制約は、\textbf{PTYレベルの行動シグナル}——インタラプトイベント、rewind操作、修正指示パターン——をファーストクラスのメモリ品質シグナルとして\textbf{欠いている}ことである。

\subsection{人間-LLMインタラクションにおける暗黙的フィードバック}

ICML 2025の論文 \cite{Liu2025ImplicitFeedback}は、WildChatおよびLMSYSデータセットにおける暗黙的ユーザーフィードバックを分析し、フィードバックの\emph{内容}が極性のみよりもモデル性能を改善することを発見した。Googleでの並行研究 \cite{Nam2025CiderChat}は、Cider Chatのテレメトリを分析し、明示的フィードバックがインタラクションのわずか0.6\%でしか発生しないことを発見した。

いずれの研究もPTYレベルのインタラプトイベントをフィードバックシグナルとして検討していない。ACMのシグナル分類体系は、ターミナルベースのコーディングエージェントに特化してこのギャップを埋める。

\subsection{検索拡張コード生成}

\textbf{Serena} \cite{OraiosAI2025Serena}は、LSPを使用してコードベースのセマンティックインデックスを構築する。ACMはSerenaと直交する。Serenaは\emph{コードベースのどの情報が利用可能か}を改善し、ACMは\emph{このコードベースでの作業に関するどの経験的パターンが利用可能か}を改善する。

\subsection{ACMの位置づけ}

表~\ref{tab:comparison}にACMと関連研究の比較を示す。

\begin{table*}[t]
\centering
\caption{ACMと関連エージェントメモリシステムの比較}
\label{tab:comparison}
\begin{tabular}{lcccc}
\toprule
\textbf{システム} & \textbf{クロスセッション} & \textbf{失敗体験} & \textbf{PTY/行動} & \textbf{コーディング} \\
 & \textbf{メモリ} &  & \textbf{シグナル} & \textbf{特化} \\
\midrule
MemGPT / Letta & Yes & No & No & No \\
Mem0 / A-MEM & Yes & No & No & No \\
A-MAC & Yes & Partial & No & No \\
REMEMBERER & Yes & Yes (Q値) & No & No \\
Self-Gen ICE & Yes & No & No & No \\
Reflexion & No* & Yes & No & No \\
ExpeL & Yes & Yes & No & No \\
Voyager & Yes & No & No & No \\
Letta Code & Yes & No & No & Yes \\
\textbf{ACM（本研究）} & \textbf{Yes} & \textbf{Yes} & \textbf{Yes} & \textbf{Yes} \\
\bottomrule
\multicolumn{5}{l}{\small *セッション内のみ。脚注参照。}
\end{tabular}
\end{table*}

\section{ACMフレームワーク}

\subsection{アーキテクチャ概要}

ACMは、既存のLLMコーディングエージェントインフラストラクチャをラップするMCP（Model Context Protocol）サーバーとして設計される。この設計選択により、Claude Code、OpenAI Codex CLI、およびMCPをサポートする任意のエージェントとの互換性が確保される。

\begin{lstlisting}[caption={ACM MCP Server hookアーキテクチャ},label={lst:hooks}]
[ACM MCP Server] -- Claude Code hooks API

  SessionStart hook
    -> ACMコンテキスト注入

  UserPromptSubmit hook
    -> インタラプト後N=3-5ターン捕捉
    -> 修正指示パターン検知

  PostToolUse hook
    -> 成功したツール完了を記録

  PostToolUseFailure hook
    -> is_interrupt=true: 失敗エントリ生成開始

  Stop hook
    -> 正常完了: 成功候補
    -> 非発火 = 補完的インタラプト検知

  SessionEnd hook
    -> エントリ確定・永続化
\end{lstlisting}

ACM層は3つの機能を実行する。

\begin{enumerate}
\item \textbf{セッション開始時の検索}: 現在のタスクに関連するエントリについて体験DBを照会し、圧縮された表現をセッションコンテキストに注入する
\item \textbf{セッション中のシグナル監視}: PTYインターフェースと会話から暗黙的フィードバックシグナルを捕捉する
\item \textbf{セッション終了時の体験書き込み}: セッションをスコアリングし、成功または失敗エントリとして保存する
\end{enumerate}

\subsection{体験エントリ構造}

ACMは2種類のエントリを管理する。

\textbf{成功エントリ}（「これをせよ」）:
\begin{lstlisting}
{
  "type": "success",
  "trigger": "Authentication bug in auth.py",
  "action": "Used FileSystemWatcher, showed
             diff before applying",
  "outcome": "Tests passed, no rewind",
  "retrieval_keys": ["auth", "JWT",
    "FileSystemWatcher", "diff"],
  "signal_strength": 0.85
}
\end{lstlisting}

\textbf{失敗エントリ}（「これをするな」）:
\begin{lstlisting}
{
  "type": "failure",
  "trigger": "Authentication bug in auth.py",
  "action": "Attempted bulk rewrite of all
             related files simultaneously",
  "outcome": "Interrupt detected. Post-
    interrupt: 'just change the diff'",
  "retrieval_keys": ["auth", "bulk rewrite"],
  "signal_strength": 0.95,
  "interrupt_context": {
    "turns_captured": 3,
    "dialogue_summary": "User wanted
      incremental diff-based changes"
  }
}
\end{lstlisting}

\subsection{検索と注入}

各セッション開始時に、ACMはユーザーの初期タスク記述からキーワードを抽出し、体験DBの\texttt{retrieval\_keys}に対してベクトル類似度検索を実行し、上位K件の関連エントリを検索して圧縮サマリーをセッションコンテキストに注入する。

注入は意図的にコンパクトである。詳細なエントリは外部に保存されパスで参照されるため、コンテキストウィンドウの膨張を回避する。これは「インデックスであってコンテンツではない」原則を実装するものである。

\subsection{体験のスコアリングと昇格}

セッション終了時に、ACMは利用可能なシグナル（第4節で詳述）を用いてセッションをスコアリングし、昇格閾値を超えるエントリを体験DBに書き込む。混合シグナルを持つセッションは、成功エントリと失敗エントリの両方を生成する場合がある。

\section{暗黙的フィードバックシグナル分類体系}

本研究の中心的な貢献は、LLMコーディングエージェントインタラクションにおいて利用可能な暗黙的フィードバックシグナルの体系的分類である。

\subsection{Level 1: インタラプト + インタラプト後の対話（最高強度）}

\textbf{シグナル}: ユーザーがエージェント実行の強制停止（Ctrl+C、Esc$\times$2）をトリガーし、その後に修正的な対話が続く。

Claude Codeの\texttt{PostToolUseFailure} hookが\texttt{is\_interrupt}ブーリアンフィールドを提供し、公式hooks APIを通じてユーザー起動のインタラプトを明示的にフラグする。

\textbf{最強である理由}: インタラプトイベントは最もコストの高いユーザーアクション——進行中のプロセスの能動的な停止——を表す。インタラプト後の対話は、ユーザーが\emph{なぜ}停止したかを明示的に表現することが多いため、特に価値がある。

\subsection{Level 2: Rewind操作（高強度）}

\textbf{シグナル}: ユーザーがrewind機能を呼び出し、会話を以前の状態に復元する。Rewindはエージェント出力の意図的な取り消しを表す——強い暗黙的拒否である。

\textbf{インタラプトとの区別}: Rewindは事後的（エージェントアクション完了後）であり、インタラプトはリアルタイム（実行中）である。

\subsection{Level 3: 修正指示（中強度）}

\textbf{シグナル}: ユーザーが会話内で自然言語の修正を提供する。言語的修正はインタラプトやrewindよりもコストが低く、軽微なミスアラインメントを反映している場合がある。ただし、セッション内での繰り返しの修正指示（3回以上）は、rewindイベントと同様に重み付けされるべきである。

\subsection{Level 4: 中断なしの完了（ベースライン / 弱い正のシグナル）}

\textbf{シグナル}: タスクがインタラプト、rewind、修正指示なしに完了する。負のシグナルの不在は成功の強い確認ではないが、テスト合格と組み合わせた場合は信頼性の高い正のシグナルである。

\subsection{シグナル強度サマリー}

注: 以下の強度スコアは実験設計のための仮説値であり、今後の実証データに基づいて較正される。序列はユーザーアクションのコストを反映する。

\begin{table}[h]
\centering
\caption{シグナル強度分類体系}
\label{tab:signals}
\small
\begin{tabular}{p{3.0cm}ccp{1.8cm}}
\toprule
\textbf{シグナル} & \textbf{スコア} & \textbf{方向} & \textbf{捕捉方法} \\
\midrule
インタラプト + 対話 & 0.90--1.00 & 負 & PTY + $N$ターン \\
Rewind & 0.75--0.90 & 負 & セッションログ \\
修正指示 (3+) & 0.60--0.80 & 負 & 会話テキスト \\
修正指示 (1) & 0.30--0.50 & 負 & 会話テキスト \\
テスト通過 + 中断なし & 0.70--0.85 & 正 & 終了コード \\
中断なし（テストなし） & 0.40--0.60 & 正 & セッションログ \\
\bottomrule
\end{tabular}
\end{table}

\section{実験設計}

\subsection{Phase 0: 実現可能性プローブ}

本格的な実験評価に先立ち、ACMの主要な技術的前提がClaude Code環境内で実現可能であるかを検証するための実現可能性プローブを実施した。

\textbf{知見1: PTYインタラプトシグナルは公式hooks APIを通じてアクセス可能である。}
Claude Codeの\texttt{PostToolUseFailure} hookは\texttt{is\_interrupt}ブーリアンフィールドを提供し、ユーザーがCtrl+Cでツール実行を中断した場合に\texttt{true}に設定される。

\textbf{知見2: セッションログは体験エントリ生成に十分なデータを提供する。}
セッションログは\texttt{\textasciitilde/.claude/projects/}配下にJSONLファイルとして保存され、完全な会話履歴、ツール呼び出し、ISO 8601タイムスタンプを含む。

\textbf{知見3: セッション後のトークン使用量は回復可能である。}
トークン使用量はセッション終了後に\texttt{\textasciitilde/.claude.json}に記録される。リアルタイムトークン数はhooks APIを通じてアクセスできない。

\textbf{知見4: Rewindイベントには専用hookが存在しない。}
ACMは\texttt{UserPromptSubmit}イベントにおける修正指示パターンとメッセージ数の減少を通じて間接的にrewindを検知する。

\subsection{リサーチクエスチョン}

\begin{description}
\item[RQ1] ACMはコンテキストウィンドウ依存を低減しつつ、タスク性能を維持・向上させるか？
\item[RQ2] ACM-S、ACM-F、ACM-SFのうち、どのメモリ構成が最大の性能向上をもたらすか？
\item[RQ3] 暗黙的フィードバックシグナルの種類は体験エントリ品質と相関するか？
\item[RQ4] ACMは人為的に制約されたコンテキストウィンドウサイズ下で持続的な性能を可能にするか？
\item[RQ5] ACMはSerenaと補完的か？組み合わせは加算的な効果をもたらすか？
\end{description}

\subsection{実験条件}

\begin{table}[h]
\centering
\caption{実験条件}
\label{tab:conditions}
\small
\begin{tabular}{ll}
\toprule
\textbf{条件} & \textbf{説明} \\
\midrule
Control & ベースLLMエージェント \\
Baseline-compact & エージェント + auto-compact \\
Baseline-Serena & エージェント + Serena \\
ACM-S & ACM（成功エントリのみ） \\
ACM-F & ACM（失敗エントリのみ） \\
ACM-SF & ACM（成功 + 失敗） \\
ACM-SF + Serena & 統合アプローチ \\
\bottomrule
\end{tabular}
\end{table}

モデル能力をアーキテクチャ効果から分離するため、固定のオープンウェイトモデルをベースエージェントとして使用する。

\subsection{コンテキストウィンドウ制約条件（RQ4用）}

\begin{table}[h]
\centering
\caption{コンテキストウィンドウ制約}
\label{tab:constraints}
\small
\begin{tabular}{ll}
\toprule
\textbf{制約} & \textbf{トークン予算} \\
\midrule
Full & 128k（モデルデフォルト） \\
Half & 64k \\
Smart zone & 50k（$\sim$40\%） \\
\bottomrule
\end{tabular}
\end{table}

\subsection{タスクスイート}

SWE-bench \cite{Jimenez2024SWEbench}はシングルセッションタスクであり、ACMのクロスセッションメモリが効果を発揮する条件を含まないため、直接使用しない。本タスクスイートはContext Rotとクロスセッション改善を示すように設計されている。

\textbf{タスクA --- マルチファイルバグ修正}: 5〜10ファイルにわたる変更を必要とする埋め込みバグを修正する。テストスイートの合格で評価。

\textbf{タスクB --- 仕様書からの機能追加}: 自然言語の仕様書に記述された機能を実装する。機能テスト + 人間レビューで評価。

\textbf{タスクC --- 設計原則に基づくリファクタリング}: 指定された設計原則をコードベース全体に一貫して適用する。自動リンティング + 人間レビューで評価。

各タスクは条件ごとに5回の反復セッションで実行される。クロスセッション改善率が主要な指標である。

\subsection{評価指標}

\begin{table}[h]
\centering
\caption{評価指標}
\label{tab:metrics}
\small
\begin{tabular}{p{2.5cm}p{2.5cm}c}
\toprule
\textbf{指標} & \textbf{測定方法} & \textbf{RQ} \\
\midrule
タスク完了率 & テスト合格 / 人間評価 & RQ1 \\
インタラプト回数 & PTY SIGINTイベント & RQ1,3 \\
Rewind回数 & セッションイベントログ & RQ1,3 \\
修正指示回数 & 会話テキスト分析 & RQ1,3 \\
コンテキスト効率 & トークン / 複雑度 & RQ4 \\
クロスセッション改善 & $\Delta$完了率 1$\rightarrow$5 & RQ1,2 \\
シグナル-品質相関 & ピアソン$r$ & RQ3 \\
\bottomrule
\end{tabular}
\end{table}

\section{議論}

\subsection{健忘エージェント問題}

現在のLLMコーディングエージェントは、設計上、健忘的である。ACMは経験的知識をセッションをまたいで持続し蓄積する形式で外部化することにより、この問題に直接対処する。

ACMを備えたエージェントは\emph{プロジェクト固有の行動パターン}を発達させることができる——コードベースに何が含まれているかだけでなく、この特定のユーザーがこの特定のコードベースで問題にどのようにアプローチすることを好むかを知ることができる。

\subsection{ファーストクラスシグナルとしてのインタラプト}

インタラプトイベントは、エージェントシステム研究においてフィードバック機構として十分に活用されていない。コーディングエージェントにおいて、インタラプトは一般的であり意味的に豊かである。誤ったアプローチを追求するエージェントは停止され、ユーザーの後続の説明は、将来そのアプローチを回避するために必要な負の訓練シグナルそのものである。

インタラプト後$N$=3〜5ターンの捕捉は、この対話がユーザーが提供する最も明示的なフィードバックであることが多いため、特に価値がある。

\subsection{強化学習との関係}

ACMのシグナル分類体系は強化学習における報酬シグナルとの構造的類似性を持つが、ACMは基盤モデルに対する重み更新を行わない——コンテキストレベルで動作し、関連する過去の体験をソフトガイダンスとして注入する。この区別により、ACMはファインチューニングインフラストラクチャを必要とせず任意のLLMでデプロイ可能となる。

\subsection{制約}

\textbf{コールドスタート}: ACMは新しいタスクタイプにおける最初のセッションでは効果を発揮しない。性能向上には蓄積された体験が必要である。

\textbf{体験の陳腐化}: コードベースの進化により、過去の体験が誤解を招く可能性がある。

\textbf{インタラプトの曖昧性}: すべてのインタラプトがエージェントのアプローチへの不満を示すわけではない——外部的な中断を反映するものもある。

\textbf{適用範囲}: 現在の設計はコーディングエージェントに特化している。他のドメインへの一般化にはシグナル分類体系の適応が必要である。

\textbf{Rewind検知}: Claude Codeはrewind操作のための専用hookイベントを提供していない。ACMは\texttt{UserPromptSubmit}イベントにおける修正指示パターンまたはメッセージ数の減少を通じて間接的にrewindを検知する。

\textbf{リアルタイムコンテキスト使用量}: アクティブセッション内のトークン数はhooks APIを通じてアクセスできない。ACMはセッション後トークン数を使用し、\texttt{PreCompact} hookをニアリミット条件のプロキシとして扱う。

\section{結論}

本研究では、暗黙的なユーザーフィードバックから成功・失敗体験を捕捉するLLMコーディングエージェントのための永続的外部メモリアーキテクチャであるAssociative Context Memory (ACM)を提示した。ACMの主要な貢献は以下の通りである。

\begin{enumerate}
\item コーディングエージェントインタラクションにおける暗黙的フィードバックシグナルの分類体系。インタラプト + インタラプト後の対話を利用可能な最強のシグナルとして同定した
\item 成功（「これをせよ」）と失敗（「これをするな」）を区別する体験エントリ構造。エージェントが正のパターンと回避行動の両方を発達させることを可能にする
\item モデル修正なしに既存のLLMコーディングエージェントエコシステムと統合するMCP互換のアーキテクチャ設計
\item クロスセッション学習とコンテキストウィンドウ依存低減を評価するための実験設計
\end{enumerate}

本研究の中心的主張——アーキテクチャ的メモリがエージェントのコンテキストウィンドウ容量への依存を低減できるということ——は、エンジニアリング実務を超えた含意を持つ。検証されれば、「健忘エージェント問題」がコンテキストウィンドウのスケーリングなしに対処可能であることを示唆し、持続的なエージェント性能への補完的な道筋を提供する。

論文草稿および関連ドキュメントは\url{https://github.com/signalcompose/okitegami_paper}で公開している。

\subsection*{著者貢献およびAI利用開示}

\textbf{知的貢献と著者責任。}
本論文におけるすべてのリサーチクエスチョン、中核的アイデア、アーキテクチャ概念、および知的貢献はHiroshi Yamatoに帰属する。著者は本論文に記載された内容に対して全責任を負う。

\textbf{研究・執筆過程におけるAIツールの使用。}
本論文は、人間主導のAI協働ワークフローを通じて制作された。

\begin{itemize}
\item \textbf{Claude.ai (Anthropic)}は\emph{Strategist}として機能した：セッション間の会話コンテキストの維持、文献検索と統合の支援、査読シミュレーション、およびテキストの起草・改訂の補助を担った。
\item \textbf{Claude Code (Anthropic)}は\emph{Implementer}として機能した：ファイル操作の実行、論文メタデータ取得のためのウェブ検索、草稿ドキュメントの更新、およびバージョン管理を担った。
\end{itemize}

AIツールが行わなかった事項：リサーチクエスチョンの着想、ACMアーキテクチャの提案、インタラプトをシグナルとして活用する着想の発見、論文の範囲や主張に関する意思決定。AIが生成したテキストはすべて、掲載前に著者によるレビュー・承認を経ている。プロセスの全容は、付随リポジトリの\texttt{docs/session-log.md}に記録されている。

\textbf{補遺 --- 2026-03-08：}
本原稿の準備過程において、第1節で述べたセッションログ維持の問題が著者自身のワークフローで観察された。本プロジェクトで使用したStrategist/Implementer協働パターン---claude.ai（Strategist）がセッションコンテキストを保持し、Claude Code（Implementer）がファイル操作を実行する---では、\texttt{docs/session-log.md}を最新に保つために手動のブリッジングが必要であった。この手動更新ステップは頻繁に漏れが生じ、セッション再開時にコンテキストのギャップが発生した。

この問題は、Claude CodeのStopフック（\texttt{.claude/hooks/update\_session\_log.py}）を実装することで対処した。このフックはClaude Codeが応答ターンを完了するたびに自動発火し、編集されたファイルの記録を\texttt{session-log.md}に追記する。フックのデバッグでは、実際のJSONLトランスクリプト構造---\texttt{tool\_use}ブロックがトップレベルではなく\texttt{assistant.message.content[]}内にネストされている---を考慮する必要があった。この一人称的観察---「記憶喪失エージェント」問題が本論文の執筆ワークフロー自体にも顕在化するという事実---は、第1節で提示したACMの動機をさらに補強するものである。

\bibliographystyle{IEEEtran}
\bibliography{../references}

\end{document}
